\subsection{线性映射的收缩的存在性}

命题 9. 设 A 为环，M 为有限生成 A-模，N 为投射 A-模，u : M → N 为 A-线性映射。

\begin{enumerate}
\def\labelenumi{\alph{enumi})}
\tightlist
\item
  设 p 为 A 的素理想。下列条件等价：
\end{enumerate}

\begin{enumerate}
\def\labelenumi{(\roman{enumi})}
\item
  存在 \(f \in \mathrm{A} - \mathfrak{p}\) 与 \(v \in \operatorname{Hom}_{\mathrm{A}_f}(\mathrm{N}_f, \mathrm{M}_f)\)，满足 \(v \circ u_f = \operatorname{Id}_{\mathrm{M}_f}\)；
\item
  存在 \(v \in \operatorname{Hom}_{\mathrm{A}_{\mathfrak{p}}}(\mathrm{N}_{\mathfrak{p}}, \mathrm{M}_{\mathfrak{p}})\)，满足 \(v \circ u_{\mathfrak{p}} = \operatorname{Id}_{\mathrm{M}_{\mathfrak{p}}}\)；
\item
  \(\kappa(\mathfrak{p})\text{-线性映射 }1 \otimes u : \kappa(\mathfrak{p}) \otimes_{\mathrm{A}} \mathrm{M} \to \kappa(\mathfrak{p}) \otimes_{\mathrm{A}} \mathrm{N}\) 是单射；
\item
  存在一个整数 \(m \geqslant 0\)、M 的元素 \(x_{1}, \ldots, x_{m}\) 以及 N 上的线性形式 \(y_{1}, \ldots, y_{m}\)，使得诸 \(x_{i}\) 在 \(\mathrm{M}_{\mathfrak{p}}\) 中的像生成 \(\mathrm{A}_{\mathfrak{p}}\)-模 \(\mathrm{M}_{\mathfrak{p}}\)，且 \(\det(<y_{j}, u(x_{i})>) \notin \mathfrak{p}\)；
\end{enumerate}

若条件 (iv) 成立，则有 \(m = [\kappa(\mathfrak{p}) \otimes_{\mathrm{A}} \mathrm{M} : \kappa(\mathfrak{p})]\)，且元素 \(1 \otimes x_i\) 构成 \(\kappa(\mathfrak{p})\)-向量空间 \(\kappa(\mathfrak{p}) \otimes_{\mathrm{A}} \mathrm{M}\) 的一个基。

\begin{enumerate}
\def\labelenumi{\alph{enumi})}
\setcounter{enumi}{1}
\tightlist
\item
  满足 a) 中条件的 A 的素理想 p 组成的集合 U 是 Spec(A) 的一个开子集，且下列条件等价：
\end{enumerate}

\begin{enumerate}
\def\labelenumi{(\roman{enumi})}
\item
  我们有 U = Spec(A)；
\item
  U 包含 A 的所有极大理想；
\item
  存在 \(v \in \operatorname{Hom}_{\mathrm{A}}(\mathrm{N}, \mathrm{M})\)，满足 \(v \circ u = \operatorname{Id}_{\mathrm{M}}\)；
\item
  u 是单射，且 Coker(u) 是投射 A-模。
\end{enumerate}

我们来证明 a)。

(i)⇒(ii)⇒(iii)：这些蕴含关系是显然的。

(iii)⇒(iv)：令 \(m = [\kappa(\mathfrak{p}) \otimes_{\mathrm{A}} \mathrm{M} : \kappa(\mathfrak{p})]\)，设 \((x_{1}, \ldots, x_{m})\) 为 M 的元素序列，使元素 \(1 \otimes x_{i}\) 构成 \(\kappa(\mathfrak{p})\)-向量空间 \(\kappa(\mathfrak{p}) \otimes_{\mathrm{A}} \mathrm{M}\) 的一个基。由 Nakayama 引理，诸 \(x_{i}\) 在 \(M_{p}\) 中的像生成 \(A_{p}\)-模 \(M_{p}\)。若条件 (iii) 成立，则元素 \(1 \otimes u(x_{i})\)（\(\kappa(\mathfrak{p})\)-向量空间 \(\kappa(\mathfrak{p}) \otimes_{\mathrm{A}} N\) 中的）线性无关。

此外，存在一个 A-模 \(N'\)、一个指标集 I，以及 A-模的同构 \(\theta : N \oplus N' \to \mathrm{A}^{(\mathrm{I})}\)，由此导出 \(\kappa(\mathfrak{p})\)-向量空间的一个同构

\[
\overline {{{\theta}}}: (\kappa (\mathfrak {p}) \otimes_ {\mathrm{A}} N) \oplus (\kappa (\mathfrak {p}) \otimes_ {A} N ^ {\prime}) \rightarrow \kappa (\mathfrak {p}) ^ {(I)}.
\]

元素 \(t_i = \bar{\theta}(1 \otimes u(x_i), 0)\)（\(\kappa(\mathfrak{p})^{(I)}\) 中的）构成一个有限自由族。于是存在 I 的元素 \(\alpha_1, \ldots, \alpha_{m}\)，使得 \(\det(\mathrm{pr}_{\alpha_j}(t_i)) \neq 0\)；N 上的线性形式 \(y_j : z \mapsto \mathrm{pr}_{\alpha_j}(\theta(z, 0))\) 即合要求。

设条件 (iv) 成立。以 \((a_{ij}) \in \mathbf{M}_m(\mathrm{A})\) 表示由系数 \(a_{ij} = \langle y_j, u(x_i)\rangle\) 组成的矩阵。设 \(g\) 为 \(\mathrm{A} - \mathfrak{p}\) 的元素，使诸 \(x_i\) 的像生成 \(\mathrm{A}_g\)-模 \(\mathrm{M}_g\) (II, § 5, 第 1 节, 命题 2)，并令 \(f = g \det(a_{ij})\)。由于 \(\det(a_{ij})\) 在 \(\mathrm{A}_f\) 中可逆，元素 \(u(x_i)\) 在 \(\mathrm{N}_f\) 中的像线性无关；因而诸 \(x_i\) 在 \(\mathrm{M}_f\) 中的像构成这个 \(\mathrm{A}_f\)-模的一个基。这就证明了 a) 的最后一个断言。现在证明 (i)。以 \(w \in \operatorname{Hom}_{\mathrm{A}}(\mathrm{N}, \mathrm{M})\) 表示映射 \(z \mapsto \sum_{j} \langle y_j, z\rangle x_j\)。我们有

\[
w \circ u (x _ {i}) = \sum_ {j} a _ {i j} x _ {j};
\]

由于诸 \(x_{i}\) 的像构成 \(M_{f}\) 的一个基，且矩阵 \((a_{ij})\) 在 \(\mathbf{M}_{m}(\mathrm{A}_{f})\) 中可逆，\((w \circ u)_{f}\)（\(M_{f}\) 的自同态）是双射的，而映射 \(v = (w \circ u)_{f}^{-1} \circ w_{f} \in \operatorname{Hom}_{\mathrm{A}_{f}}(\mathrm{N}_{f}, \mathrm{M}_{f})\) 满足条件 (i)。

我们来证明 b)。U 是开集这一事实由 a) 的条件 (i) 得出。

(iii)⇒(i)⇒(ii)：这是显然的。

(iv)⇒(iii)：在 (iv) 的假设下，序列 0 → M \(\xrightarrow{u}\) N \(\longrightarrow\) Coker(u) → 0 是正合且分裂的，故得 (iii)。

\begin{enumerate}
\def\labelenumi{(\roman{enumi})}
\setcounter{enumi}{1}
\tightlist
\item
  (ii) \(\Rightarrow\) (iv)：如上引入一个 A-模同构 \(\theta: \mathrm{N} \oplus \mathrm{N}' \to \mathrm{A}^{(\mathrm{I})}\)。设 \(u'\) 为从 M 到 \(\mathrm{A}^{(\mathrm{I})}\) 中的映射（由 \(u'(x) = \theta(u(x), 0)\) 定义的）。存在 I 的一个有限子集 J，使 \(u'\) 的像含于子模 \(\mathrm{A}^{\mathrm{J}}\)（\(\mathrm{A}^{(\mathrm{I})}\) 的）。以 \(u'': \mathrm{M} \to \mathrm{A}^{\mathrm{J}}\) 表示由 \(u'\) 诱导的映射。在假设 (ii) 之下，对 A 的每个极大理想 m，\(\mathrm{A}_{\mathfrak{m}}\)-线性映射 \(u_{\mathfrak{m}}'\)（从 \(\mathrm{M}_{\mathfrak{m}}\) 到 \(\mathrm{A}_{\mathfrak{m}}^{(\mathrm{I})}\) 中的）有一个收缩，因而 \(u_{\mathfrak{m}}''\) 也有；于是 \(u_{\mathfrak{m}}''\) 是单射，且它的像是 \(\mathrm{A}_{\mathfrak{m}}^{\mathrm{J}}\) 中的一个直和项，从而它的余核是投射 \(\mathrm{A}_{\mathfrak{m}}\)-模。A-模 \(\operatorname{Coker}(u'')\) 按构造是有限表现的；因而是投射的 (II, § 5, 第 2 节, 定理 1)。同态 \(u''\) 是单射 (II, § 3, 第 3 节, 定理 1)；从而 \(u\) 是单射。A-模 \(\operatorname{Coker}(u')\) 一方面同构于 \(\operatorname{Coker}(u) \oplus \mathrm{N}'\)，另一方面同构于 \(\operatorname{Coker}(u'') \oplus \mathrm{A}^{(\mathrm{I}-\mathrm{J})}\)。由于 A-模 \(\mathrm{A}^{(\mathrm{I}-\mathrm{J})}\)、\(\operatorname{Coker}(u'')\) 与 \(\mathrm{N}'\) 都是投射的，可知 \(\operatorname{Coker}(u)\) 是投射的，这就完成了 (iv) 的证明。
\end{enumerate}
% semantic-fragments: legacy-input-seam
\relax{}
